\begin{recipe}
[% 
    preparationtime = {\unit[30]{min}},
    bakingtime={2 x \unit[8--12]{min}},
    bakingtemperature={\protect\bakingtemperature{
        topbottomheat=\unit[250]{\textcelcius}}},
    portion = {\portion{2}},
    source = {Adrian Hieber},
    rating = {3},
    calory = {915 kcal | 44g Protein | 95g KH | 40g Fett},
    price = {1,77€}
]
{Flammkuchen}
    
\graph{
    big=pic/hauptgericht/04_flammkuchen
}

\introduction{%
    Selbstgemachter Flammkuchen geht einfacher als man denkt. 
    Dünner Teig, cremiger Belag und richtig heiß gebacken -- 
    perfekt für einen entspannten Abend oder wenn Besuch da ist.
    
}

\ingredients{%
    215 g & Mehl\index{Getreide \& Teigwaren!Mehl}\\
    1{,}5 g & Trockenhefe\\
    5 g & Salz\\
    1{,}5 g & Zucker\\
    10 g & Olivenöl\\
    130 g & Wasser\\
    125 g & Creme (z.B. vegane Creme oder Creme fraîche)\index{Dairy!Creme Fraiche}\\
    50 g & Joghurt oder Sahne\index{Dairy!Sahne}\\
    \nicefrac{1}{2} & Rote Zwiebel\index{Gemüse!Zwiebel}\\
    175 g & Räuchertofu\index{Tofu \& Fleischalternativen!Tofu}\\
    1 Msp. & Muskat\\
    1 EL & Schnittlauch
}

\preparation{% 
    \step%
    \textbf{Teig:} Mehl, Hefe, Salz und Zucker vermischen. 
    Wasser und Olivenöl hinzufügen und zu einem geschmeidigen Teig kneten.
        
    \step%
    Teig halbieren und auf Backpapier sehr dünn ausrollen.
        
    \step%
    \textbf{Belag:} Zwiebel in feine Streifen schneiden. 
    Räuchertofu in kleine Würfel schneiden.
        
    \step%
    Creme, Joghurt/Sahne, Schnittlauch und Muskat verrühren.
        
    \step%
    Flammkuchen dünn mit der Creme bestreichen und mit Tofu und Zwiebeln belegen.
        
    \step%
    Bei \unit[250]{\textcelcius} etwa 8--12 Minuten backen, bis der Rand goldbraun ist. Vorgang mit dem zweiten Flammkuchen wiederholen.
}

\hint{%
    Je dünner der Teig, desto besser wird der Flammkuchen. 
    Der Ofen sollte wirklich gut vorgeheizt sein.
}

\end{recipe}